\begingroup\small\setlength{\tabcolsep}{3pt}
\begin{longtable}{lrrr}
\caption{Sensibilidad del gradiente patrimonial de incorporación}\label{tab:calidad-tesis_calidad_gradiente}\\
\toprule
Vía/ventana & Edad3/celda20 pp & Edad2/celda30 pp & N soporte Q1/Q4 \\
\midrule\endfirsthead
\multicolumn{4}{l}{Continuación del cuadro \thetable}\\
\toprule
Vía/ventana & Edad3/celda20 pp & Edad2/celda30 pp & N soporte Q1/Q4 \\
\midrule\endhead
\bottomrule\endfoot
K 2001--2003 & -0.5 [-3.1; 1.9] & -0.8 [-3.2; 1.7] & 956/831 \\
K 2012--2015 & 0.5 [-4.0; 5.0] & 0.3 [-4.2; 4.9] & 918/894 \\
K 2018--2021 & -4.0 [-7.7; -0.4] & -4.0 [-7.6; -0.4] & 668/596 \\
A 2003--2012 & 6.8 [0.3; 13.2] & 6.9 [0.9; 13.3] & 527/499 \\
A 2012--2015 & 4.0 [-0.8; 8.6] & 4.0 [-1.0; 8.8] & 919/895 \\
A 2015--2018 & 6.2 [0.7; 11.6] & 4.7 [-0.8; 10.2] & 508/490 \\
\end{longtable}
{\footnotesize Fuente: auditoría10, insumos01--09 intactos. Muestra sin ponderación; IC condicionales, sin inferencia nacional o causal. P: primera vía observada; C: otra vía basal afirmativa. Q1 menos Q4 entre negativas K/A reobservadas; sexo y escolaridad se mantienen. Edad2=60--69/70+, edad3=60--64/65--69/70+. Cambian soporte y mezcla, no es la misma población. Todas las cuatro especificaciones están en el informe; no se elige por significancia.\par}
\endgroup
